\chapter{Two Paths to Conformance Testing: Fuzzing and Symbolic Execution}\label{chapter:paper4}
Network protocol conformance can be assessed through different testing techniques, each offering distinct strengths and limitations.
The previous chapters have focused on symbolic execution (SE) as a principled technique for systematically exploring protocol behavior and for reasoning about requirement violations that arise only under narrowly constrained conditions.
By progressively extending SE with requirement-driven oracles, external monitors, and differential analysis, we demonstrated how it can be applied effectively to complex, stateful protocols such as DTLS and QUIC.

However, SE is not the only technique used in practice to test protocol implementations.
Coverage-guided fuzzing has become a widely adopted approach for testing networked systems, particularly in settings where scalability and low analysis overhead are important.
Fuzzers excel at exercising large input spaces quickly and can expose robustness issues as well as logical flaws that manifest through externally observable behavior.
At the same time, fuzzing is not typically employed to check protocol requirements or to systematically target specific semantic conditions.

This chapter applies the requirement-centric perspective developed in earlier chapters to both techniques and compares their effectiveness when applied to the same requirements and implementations.
In particular, it examines how fuzzing and SE differ in their ability to expose requirement violations in the context of an IoT protocol.

\section{Technique}
\subsection{Test Harness Creation}
Testing protocol conformance requires a controlled execution environment in which interactions with a protocol entity can be reproduced and manipulated systematically.
As in the previous chapters, this is achieved through a dedicated test harness that mediates communication between the system under test (SUT) and its peer and that provides full control over the messages exchanged during a protocol interaction.

The construction of the harness begins by capturing a representative interaction with the protocol entity.
For CoAP, this interaction consists of a sequence of request and response messages exchanged between a client and a server.
Packet capture is performed by configuring sample programs provided with the SUT to execute a valid interaction and recording the packets sent to the protocol entity using standard network analysis tools such as TCPdump or Wireshark.
The resulting packet sequence serves as a concrete reference interaction that can be replayed during testing.

During testing, the harness advances the interaction step by step by replaying the pre-captured packets.
At each step, the harness invokes the appropriate API of the SUT to deliver an incoming message and to obtain the corresponding response.
Responses produced by the SUT are recorded and used by the harness to determine how the interaction proceeds.
This execution model ensures that the protocol logic exercised during testing follows the same high-level interaction pattern as the captured execution, while allowing individual message fields to be varied or constrained as required by the testing technique.

The test harness should be designed to support all protocol requirements considered for testing.
Depending on the nature of the requirements, this may involve executing a single interaction that exercises multiple protocol features or running several distinct interactions that target different parts of the protocol.
In either case, the harness provides a uniform mechanism for delivering inputs, observing outputs, and controlling progress through the protocol interaction.

A key design goal of the harness is to enable both fuzzing and SE with minimal adaptation.
The same harness structure is used in both settings, with differences limited to how input messages are generated and how responses are analyzed.
For fuzzing, the harness supplies mutated variants of the captured packets, while for SE it designates selected fields as symbolic and relies on the symbolic executor to explore feasible variations.
By sharing a common harness, both techniques are applied to identical interaction patterns and protocol requirements, enabling a meaningful comparison of their effectiveness in exposing requirement violations.

\subsection{Encoding Protocol Requirements as Assertions}
Network protocol specifications define requirements not only over individual messages but also over sequences of messages exchanged between protocol entities.
For CoAP, such requirements are specified across several RFC documents, most notably RFC~7252~\cite{RFC7252} and RFC~7959~\cite{RFC7959}.
Once relevant requirements have been extracted from these specifications, they can be translated into executable assertions that are evaluated during a protocol interaction.

In this work, protocol requirements are encoded as assertions over the packets exchanged during an interaction.
These assertions are inserted at points in the SUT where response messages are generated, allowing the observed behavior of the implementation to be checked against the specification at runtime.
In the CoAP implementations considered, assertions were inserted into the functions responsible for transmitting response messages.
These functions already operate on structured representations of outgoing packets, which allows assertions to directly inspect response fields.
Fields of the corresponding incoming packet are accessed through a global pointer that is maintained by the test harness and updated at each interaction step.

We illustrate this approach by presenting the encoding of two representative CoAP requirements.


\paragraph{Message ID Echo.}
The CoAP specification mandates that the Message ID of certain request messages must be echoed in the corresponding response.
RFC~7252~\cite[p.~24]{RFC7252} states:
%
\rfcdisplayquote{%
	The Message ID is a 16-bit unsigned integer that is generated by the sender of a Confirmable or Non-confirmable message and included in the CoAP header.
	The Message ID MUST be echoed in the Acknowledgement or Reset message by the recipient.}%
%
Let $\rfcm{p_{in}}$ denote the parsed incoming packet and $\rfcm{p_{out}}$ the response packet generated by the SUT.
The requirement can be encoded by the following assertion, which is evaluated when the response packet is constructed:
%
\begin{equation*}
	\begin{split}
		& \scriptstyle \mathtt{assert} \, (((\rfcm{p_{in}.type}\, =\, \rfcm{CON} \, \lor \, \rfcm{p_{in}.type}\, =\, \rfcm{NON}) \, \land \, \\[-5pt]
		& \scriptstyle \quad\quad \,\,\,\,\,\,\,(\rfcm{p_{out}.type}\, =\, \rfcm{ACK} \, \lor \, \rfcm{p_{out}.type}\, =\, \rfcm{RST})) \implies \\[-5pt]
		& \scriptstyle \qquad\qquad\qquad \rfcm{p_{in}.message\_id}\, =\, \rfcm{p_{out}.message\_id})
	\end{split}
\end{equation*}
%
This assertion ensures that whenever the SUT generates an acknowledgement or reset in response to a Confirmable or Non-confirmable request, the Message ID is correctly echoed.
During testing, the assertion is evaluated at each interaction step, and any violation immediately indicates nonconformance with the specification.

\paragraph{Further Request Block Size.}
Some CoAP requirements relate the handling of a request to information conveyed in earlier responses.
RFC~7959~\cite[p.~14]{RFC7959} specifies such a relationship for block-wise transfers:
%
\rfcdisplayquote{%
	In response to a request with a payload (e.g., a PUT or POST transfer), the block size given in the Block1 Option indicates the
	block size preference of the server for this resource.
	\dots
	the client SHOULD heed the preference indicated and, for all further
	blocks, use the block size preferred by the server or a smaller one.}%
%

To encode this requirement, the assertion logic explicitly maintains state across packets.
A boolean variable $\rfcm{first\_response}$ indicates whether the current response is the first generated by the server, and a variable $\rfcm{preferred\_size}$ stores the block size advertised by the server.
The requirement is checked using the following logic:
%
\begin{equation*}
	\begin{split}
		& \scriptstyle \mathtt{if} \, (\, \rfcm{first\_response} \, \land \, \rfcm{p_{in}.code} \, \in \, \rfcm{REQ\_CODES})\\[-5pt]
		& \scriptstyle \quad \rfcm{preferred\_size} \, := \, \rfcm{BlockSize (p_{out})}\\[-5pt]
		& \scriptstyle \quad \rfcm{first\_response} \, := \, \mathtt{false} \\[-5pt]
		& \scriptstyle \mathtt{else \, if} (\neg \, \rfcm{first\_response})\\[-5pt]
		& \scriptstyle \quad \mathtt{assert} ((\rfcm{BlockSize (p_{in})} \, > \, \rfcm{preferred\_size}) \, \implies \\[-5pt]
		& \scriptstyle \qquad\qquad\qquad \rfcm{p_{out}.code} \, \in \,  \rfcm{ERR\_CODES})
	\end{split}
\end{equation*}
%
The assertion is evaluated only after the server’s preferred block size has been established by the first relevant response.

Here, \rfc{REQ\_CODES} represents the set of CoAP request codes.
The function $\rfcm{BlockSize}$ extracts the block size specified by the Block1 option of a packet, and \rfc{ERR\_CODES} denotes the set of CoAP error response codes.
The assertion checks that if a client continues to send request blocks larger than the server's advertised preference, the server responds with an error.

While the RFC does not prescribe a specific error response for this case, requiring an error response provides a reasonable approximation of the intended behavior and allows violations to be detected systematically.

Once the SUT has been extended with such assertions, they serve as executable representations of protocol requirements.
Both fuzzing and SE can then be applied to explore whether the implementation can be driven into executions that violate these assertions, thereby revealing instances of protocol nonconformance.

\subsection{Testing Requirements Using Fuzzing}
While coverage-guided fuzzing is traditionally used to discover crashes and robustness issues, it can be adapted to support requirement-driven testing when combined with explicit assertions.
In this work, we use fuzzing not as a robustness testing technique, but as a mechanism for exploring executions that may violate protocol requirements encoded as runtime assertions.
Standard fuzzing infrastructure can be leveraged for this purpose without modifying the fuzzer itself, by generating inputs that may trigger violations of the encoded assertions.
In our evaluation, we use AFL++~\cite{AFLplusplus-Woot20}, a state-of-the-art fuzzer that incorporates techniques from recent research and provides efficient feedback-driven input generation.

Fuzzing is performed by executing the test harness described earlier.
As initial seeds, we use valid, pre-captured packets obtained from protocol interactions executed using the sample programs provided with each implementation.
These packets serve as well-formed starting points from which the fuzzer generates mutated inputs.
During a fuzzing campaign, the fuzzer repeatedly executes the harness and supplies mutated variations of the incoming packets to the SUT.
All encoded assertions are evaluated during each execution, and any violation is immediately detected.

Some protocol requirements, particularly those related to block-wise transfer in CoAP, depend on the processing of multiple packets and therefore require that protocol state be built across several interaction steps before an assertion can be checked.
Standard fuzzers are not designed to reason about such stateful interactions.
To accommodate these requirements, we represent the packets of a complete interaction as a single contiguous input buffer, with each packet placed at a fixed offset.
During execution, the harness delivers packets to the SUT by reading from the appropriate offsets within this buffer.
Since the fuzzer is free to arbitrarily mutate the input buffer, packet boundaries are not guaranteed to remain intact.
As a result, some mutated inputs may cause packet boundaries to shift or packets to become malformed.
In such cases, the harness may fail to parse the input correctly or the SUT may immediately reject the input, and these executions do not contribute to requirement checking.
When packet parsing succeeds, however, this design allows the fuzzer to mutate a single input buffer while still exercising a multi-packet protocol interaction.
As a result, state-dependent requirements can be tested using standard fuzzing infrastructure without modifying the fuzzer itself.

Taking into account the inherent randomness of fuzzing, campaigns are repeated multiple times to increase confidence in their ability to expose violations.
When the fuzzer triggers an assertion violation, it stores the corresponding mutated input together with metadata such as the time of discovery and the number of mutations applied.
These concrete test cases can then be analyzed to determine which packet fields and values caused the requirement violation and to reproduce the behavior on the unmodified implementation.

Overall, fuzzing provides a lightweight and scalable means of exploring large portions of the input space and of exercising multiple requirements simultaneously.
However, its effectiveness depends on the fuzzer's ability to reach inputs that satisfy the specific conditions under which assertions fail.
In the next subsection, we contrast this behavior with SE, which reasons explicitly about such conditions and can systematically target individual requirements.

\subsection{Testing Requirements Using Symbolic Execution}
The approach for testing requirements using SE employed in this work closely follows the requirement-driven SE technique introduced in Chapter~\ref{chapter:paper1}, using the same test harness structure and execution model.
The main differences in this work concern the protocol under test and the specific requirements considered, rather than the underlying SE methodology.
To test protocol requirements using SE, we employ KLEE~\cite{KLEE@OSDI-08}, which we use as the SE engine throughout this thesis.
As in the fuzzing-based approach, SE is performed by executing the test harness described earlier.
However, unlike fuzzing, SE targets individual requirements explicitly and is therefore performed in separate runs, with each run checking exactly one requirement.

In each SE run, KLEE executes the test harness using a pre-captured sequence of valid packets that define the interaction under test.
Before a packet is processed by the protocol party, the fields relevant to the requirement being checked are designated as symbolic, while all other fields remain concrete.
This selective symbolic treatment restricts exploration to the part of the input space that is relevant for the requirement and avoids unnecessary path explosion.

As in the earlier chapters, packets are parsed into structured representations to allow symbolic manipulation of individual protocol fields, and are serialized back into byte buffers before being processed by the protocol implementation.

As an example, to test the Message ID Echo requirement, the \rfc{message\_id} and \rfc{type} fields of the incoming packet are made symbolic.
KLEE then explores all feasible execution paths that arise from different assignments to these fields.
If there exists an execution path along which the corresponding assertion is violated, KLEE reports the violation and produces a concrete assignment for the symbolic fields that triggers it.
This assignment constitutes a concrete test case that demonstrates nonconformance with the requirement.

In this setting, a test case corresponds to the specific values assigned to the symbolic protocol fields, rather than to an entire packet sequence.
These values can be substituted back into the original packet structure and replayed against the unmodified implementation to reproduce the violation.
Symbolic execution therefore provides a systematic means of determining whether a requirement can be violated at all and, if so, of generating precise witnesses that expose the violation.

\section{Evaluation}
The purpose of the evaluation is to assess how fuzzing and SE behave in practice when applied to the same protocol requirements within a shared testing framework.
The evaluation focused on a set of representative protocol requirements extracted from the CoAP specification and its extensions.

Overall, our findings demonstrate that both fuzzing and SE are effective at exposing protocol nonconformances when requirements are encoded explicitly as assertions.
Across the evaluated requirements, both fuzzing and SE identified the same set of nine distinct nonconformances in the tested implementations.
All identified issues were reported to the respective developers, and most have been fixed at the time of writing.
These results indicate that assertion-based testing is capable of uncovering concrete and actionable specification violations in real-world CoAP implementations.

In terms of effectiveness at detecting violations, we did not observe a substantial difference between the two techniques.
For the evaluated requirements, fuzzing and SE were able to identify violations for the same subsets of requirements in each implementation.
Moreover, violations were typically discovered quickly.
Fuzzing consistently generated inputs that triggered assertion failures within seconds for most requirements, while SE also produced violating inputs within comparable time frames in the majority of cases.
This suggests that, when requirement violations are reachable through relatively shallow protocol interactions, both techniques are well suited to exposing them.

The main distinction between fuzzing and SE becomes apparent when considering requirements that an implementation respects.
Fuzzing is inherently probabilistic and cannot provide evidence that a requirement is upheld beyond the absence of observed violations within a given time budget.
Symbolic execution, by contrast, explores program paths systematically and can, in principle, establish that no execution violating a requirement exists within the explored search space.
In practice, this ability is limited by path explosion and solver performance.
For some requirements, SE was able to complete exploration and thus provide stronger assurance of correctness.
For others, exploration did not complete within the imposed time limits.

Taken together, these observations suggest a complementary relationship between the two techniques.
Fuzzing offers a lightweight and scalable way to uncover specification violations quickly, making it well suited for early testing.
Symbolic execution, while more computationally expensive, provides deeper insight into the conditions under which violations occur and can offer stronger guarantees when exploration completes.
Our experience indicates that combining both techniques within a single assertion-based framework allows developers to benefit from the strengths of each while mitigating their respective limitations.
A detailed description of the experimental setup, the evaluated implementations, and the complete set of discovered issues is provided in Paper \URoman{4}.

